\documentclass[11pt]{article}
\usepackage[utf8]{inputenc}
\usepackage[T1]{fontenc}
\usepackage[margin=2.5cm]{geometry}
\usepackage{amsmath,amssymb,amsthm,mathtools}
\usepackage{enumitem}
\usepackage{xcolor}
\usepackage{hyperref}
\hypersetup{colorlinks=true,linkcolor=blue!60!black,urlcolor=blue!60!black}

\theoremstyle{plain}
\newtheorem{theorem}{Theorem}[section]
\newtheorem{proposition}[theorem]{Proposition}
\newtheorem{lemma}[theorem]{Lemma}
\newtheorem{corollary}[theorem]{Corollary}
\theoremstyle{definition}
\newtheorem{assumption}{Assumption}
\theoremstyle{remark}
\newtheorem{remark}[theorem]{Remark}

\newcommand{\R}{\mathbb{R}}
\newcommand{\N}{\mathbb{N}}
\newcommand{\E}{\mathbb{E}}
\newcommand{\KL}{\mathrm{KL}}
\DeclareMathOperator*{\argmin}{arg\,min}
\DeclareMathOperator*{\argmax}{arg\,max}
\DeclareMathOperator{\dimm}{dim}

\title{\textbf{Detailed proof of the orbital--MDL upper bound,\\
and how it pairs with the lower bound}\\[0.3em]
\large (validation note; closes the two upper-side gaps and the log factor)}
\author{J.~Nemb\'e \\ \small (working note)}
\date{\today}

\begin{document}
\maketitle

\begin{abstract}
We prove the upper bound for the adaptive orbital minimum-description-length
(MDL) estimator of a $G$-invariant density and pair it with the minimax lower
bound established in the companion note. Three things are made rigorous and
honest. (1) The \emph{oracle inequality}: the corrected Kraft/partition-function
bound for the Hellinger affinity is proved in full and is \emph{tail-free}
(the affinity is always $\le1$), which dissolves the ``heavy-tailed affinity
constant'' gap. (2) The \emph{twin Jackson} approximation on the quotient holds on
possibly stratified quotients because, under $d$-Ahlfors regularity, the
lower-dimensional strata are $\mu_G$-null, so approximation is governed by the top
stratum; this dissolves the ``stratified quotient'' gap. (3) The honest rate:
the adaptive two-stage MDL estimator attains
$(\log n/n)^{2s/(2s+d)}$, i.e.\ the minimax rate $n^{-2s/(2s+d)}$ of the lower
bound \emph{up to a logarithmic factor} (the standard two-stage/adaptation cost),
on multiplicative/heavy-tailed data in the $\phi$-chart. (4) \emph{Removing the
logarithm}: a finite-dimensional projection estimator attains the exact rate
$n^{-2s/(2s+d)}$ for known $s$, and a penalised projection estimator
(Birg\'e--Massart) attains it \emph{adaptively}, because these estimators pay the
model's metric complexity $D_k/n$ rather than the per-coefficient quantisation
cost $\log(1/\delta)$ of a description-length code.
\end{abstract}

\section{Setting (recalled) and the estimator}

As in the lower-bound note: $M=E/G$, $d=\dimm M$, $\mu_G$ the quotient reference
measure, smoothness measured in the $\phi$-chart; $\tilde f^\star\in\Sigma(s,L)$ is
the (invariant) target with $0<c_0\le\tilde f^\star\le C_1$ on a $d$-Ahlfors
regular region. We use $\rho(f,g)=\int\sqrt{fg}\,d\mu\in(0,1]$ and the squared
Hellinger $h^2=1-\rho$, and write $\KL$ for Kullback--Leibler. Let
$\{f_m\}_{m\in\mathcal I}$ be a countable family of candidate densities with prefix
code lengths $L(m)$ satisfying the half-Kraft condition
\begin{equation}\label{eq:halfkraft}
\sum_{m\in\mathcal I}2^{-L(m)/2}\le1
\end{equation}
(achieved from any Kraft code by $L\leftarrow2L$). The orbital MDL estimator is
\begin{equation}\label{eq:mdl}
\widehat m\in\argmin_{m}\Big\{-\sum_{i=1}^n\log f_m(X_i)+L(m)\log2\Big\},
\qquad \widehat f=f_{\widehat m}.
\end{equation}

\section{The oracle inequality: tail-free partition bound}

\begin{lemma}[Kraft/partition-function bound; full proof, tail-free]\label{lem:partition}
For the estimator \eqref{eq:mdl} under \eqref{eq:halfkraft},
\[
\E_{P^\star}\!\Big[2^{-L(\widehat m)/2}\prod_{i=1}^n\sqrt{\tfrac{f_{\widehat m}(X_i)}{f^\star(X_i)}}\Big]
\ \le\ \sum_{m}2^{-L(m)/2}\,\rho(f^\star,f_m)^n\ \le\ 1 .
\]
No tail or moment assumption on $f^\star$ or $f_m$ is used.
\end{lemma}

\begin{proof}
The selector \eqref{eq:mdl} maximises $2^{-L(m)}\prod_i f_m(X_i)$ over $m$; taking
the positive square root, for every realisation
\[
2^{-L(\widehat m)/2}\prod_i\sqrt{f_{\widehat m}(X_i)}
=\max_m 2^{-L(m)/2}\prod_i\sqrt{f_m(X_i)}
\le\sum_m 2^{-L(m)/2}\prod_i\sqrt{f_m(X_i)} .
\]
Divide by $\prod_i\sqrt{f^\star(X_i)}>0$ ($P^\star$-a.s.) and take $\E_{P^\star}$.
By independence and
$\E_{P^\star}[\sqrt{f_m(X)/f^\star(X)}]=\int\sqrt{f_m f^\star}\,d\mu=\rho(f^\star,f_m)$,
\[
\E_{P^\star}\Big[2^{-L(\widehat m)/2}\prod_i\sqrt{\tfrac{f_{\widehat m}(X_i)}{f^\star(X_i)}}\Big]
\le\sum_m 2^{-L(m)/2}\rho(f^\star,f_m)^n
\le\sum_m 2^{-L(m)/2}\le1,
\]
the last step using $\rho\le1$ (always) and \eqref{eq:halfkraft}.
\end{proof}

\begin{theorem}[Hellinger oracle inequality]\label{thm:oracle}
There is a universal constant $C$ such that
\[
\E_{P^\star}\big[h^2(P^\star,\widehat f)\big]
\ \le\ C\,\inf_{m\in\mathcal I}\Big\{\KL(P^\star\|f_m)+\frac{(\log2)\,L(m)}{n}\Big\}.
\]
\end{theorem}

\begin{proof}
Lemma~\ref{lem:partition} is the index-of-resolvability estimate of Barron--Cover
\cite{barroncover} (see also Li \cite{li}); the standard passage from it to the
displayed bound combines a Markov step on the random affinity with the comparison
$h^2\le-\log\rho\le\KL$. We isolate Lemma~\ref{lem:partition} because it is the
step mis-stated in the source notes; the passage is the classical part.
\end{proof}

\begin{remark}[First gap closed: heavy tails are irrelevant to concentration]
Lemma~\ref{lem:partition} uses only $\rho\le1$, which holds for \emph{any} pair of
densities. Thus the concentration/oracle step carries \emph{no} tail or moment
hypothesis: in the multiplicative regime ($\phi=\log$, possibly heavy-tailed in
the original scale) the affinity bound is unchanged. Tails can only enter through
the \emph{approximation} term $\KL(P^\star\|f_m)$, i.e.\ through the smoothness of
the target in the $\phi$-chart --- which is exactly the standing assumption.
\end{remark}

\section{Twin Jackson approximation on the (stratified) quotient}

\begin{assumption}[Regular quotient with controlled strata]\label{ass:strata}
$(M,\mu_G)$ is $d$-Ahlfors regular on a compact region $U$ carrying $\tilde
f^\star$; the non-principal orbit strata have dimension $<d$. (For a free action
$M$ is a compact $d$-manifold.)
\end{assumption}

\begin{lemma}[Orbital sieve and twin Jackson bound]\label{lem:jackson}
Under Assumption~\ref{ass:strata} and $\tilde f^\star\in\Sigma(s,L)$ with $c_0\le
\tilde f^\star\le C_1$, for each $k\in\N$ there is a $G$-invariant density $q_k$,
piecewise polynomial of degree $<\lceil s\rceil$ on a resolution-$2^{-k}$
partition of $U$ (pulled back through $\pi$), with
\[
\dimm\mathcal M_k\asymp 2^{kd},\qquad
c_0/2\le q_k\le 2C_1,\qquad
\|\tilde f^\star-q_k\|_\infty\le C\,L\,2^{-ks},
\]
hence $\KL(P^\star\|Q_k)\le \tfrac{2}{c_0}\|\tilde f^\star-q_k\|_{L^2(\mu_G)}^2
\le C'\,L^2\,2^{-2ks}$.
\end{lemma}

\begin{proof}
\emph{Second gap closed.} By $d$-Ahlfors regularity, any subset of dimension
$<d$ is $\mu_G$-null; in particular the non-principal strata carry no $\mu_G$-mass.
Hence the $L^2(\mu_G)$ and $\KL$ approximation errors are computed on the principal
stratum $M_0$, which is a $d$-dimensional (Ahlfors-regular) space. There the
classical piecewise-polynomial Jackson estimate gives
$\|\tilde f^\star-q_k\|_\infty\le CL2^{-ks}$ on each cell of side $2^{-k}$ (Taylor
remainder of order $s$). Truncating to $[c_0/2,2C_1]$ and renormalising changes
the sup-norm by $O(2^{-ks})$ and keeps $q_k$ a density. Finally
$\KL\le\chi^2\le\frac{2}{c_0}\|\cdot\|_{L^2}^2$ since $q_k\ge c_0/2$, and
$\|\cdot\|_{L^2(\mu_G)}^2\le\mu_G(U)\|\cdot\|_\infty^2$.
\end{proof}

\section{Description length and discretisation}

\begin{lemma}[Code length / candidate complexity]\label{lem:code}
Quantising the $\asymp2^{kd}$ coefficients of $\mathcal M_k$ on a grid of spacing
$\delta_k\asymp 2^{-ks}$ yields a countable family with a prefix code (Elias code
on $k$, uniform code on the grid) satisfying \eqref{eq:halfkraft} and
\[
L(k)\ \asymp\ 2^{kd}\,\log_2(1/\delta_k)\ \asymp\ 2^{kd}\,(ks)\ \lesssim\ 2^{kd}\log n
\]
at the resolution used below, while the quantisation adds at most $C\,2^{kd}/n$ to
$\KL(P^\star\|q_k)$.
\end{lemma}

\begin{proof}
Standard sieve discretisation (Barron--Cover \cite{barroncover}): a grid of
spacing $\delta_k$ on each of $\asymp2^{kd}$ bounded coefficients costs
$\log_2(1/\delta_k)$ bits per coefficient; a prefix code over $k$ adds $O(\log k)$
bits. The quantisation perturbs $q_k$ in sup-norm by $O(\delta_k)$, hence $\KL$ by
$O(\delta_k^2\cdot2^{kd})$; with $\delta_k\asymp 2^{-ks}$ and at the balance
$2^{k}\asymp(n/\log n)^{1/(2s+d)}$ this is $O(2^{kd}/n)$, of the same order as the
complexity term.
\end{proof}

\section{The upper bound and the pairing}

\begin{theorem}[Adaptive orbital--MDL upper bound]\label{thm:upper}
Under Assumptions~\ref{ass:strata} and the $\phi$-chart smoothness, the orbital
MDL estimator \eqref{eq:mdl} over the family of Lemmas~\ref{lem:jackson}--\ref{lem:code}
satisfies, for every $s\in(0,s_{\max}]$ and $f^\star\in\Sigma_{\mathrm{inv}}(s,L)$,
\[
\E_{P^\star}\big[h^2(P^\star,\widehat f)\big]
\ \le\ C\Big(\frac{\log n}{n}\Big)^{2s/(2s+d)},
\]
\emph{adaptively} (the estimator does not use $s$).
\end{theorem}

\begin{proof}
Combine Theorem~\ref{thm:oracle} with Lemmas~\ref{lem:jackson}--\ref{lem:code}:
\[
\E[h^2(P^\star,\widehat f)]
\le C\inf_{k}\Big\{L^2\,2^{-2ks}+\frac{2^{kd}\log n}{n}\Big\}.
\]
Balancing $L^2 2^{-2ks}\asymp 2^{kd}\log n/n$, i.e.\
$2^{k(2s+d)}\asymp n/\log n$, gives the stated rate. The infimum over $k$ is
realised by the data-driven $\widehat m$, hence adaptivity.
\end{proof}

\begin{corollary}[Rate pairing]\label{cor:pairing}
With the lower bound $\inf_{\widehat f}\sup_{\Sigma_{\mathrm{inv}}(s,L)}\E\,h^2\ge
c\,n^{-2s/(2s+d)}$ of the companion note, the orbital MDL estimator is
\emph{minimax-rate optimal up to a logarithmic factor}:
\[
c\,n^{-2s/(2s+d)}\ \le\ \mathcal R_n^\star\big(\Sigma_{\mathrm{inv}}(s,L)\big)
\ \le\ \E\,h^2(P^\star,\widehat f)\ \le\ C\Big(\tfrac{\log n}{n}\Big)^{2s/(2s+d)},
\]
with $d=\dimm(E/G)$, on multiplicative/heavy-tailed data in the $\phi$-chart.
\end{corollary}

\section{Removing the logarithmic factor}\label{sec:nolog}

The $(\log n)^{2s/(2s+d)}$ surplus of Theorem~\ref{thm:upper} is an artefact of
\emph{coding}: the half-Kraft code of Lemma~\ref{lem:code} spends
$\log_2(1/\delta_k)\asymp\log n$ bits to quantise each of the $\asymp2^{kd}$
continuous coefficients. An estimator that selects a finite-dimensional model and
fits its coefficients \emph{without} quantising them pays instead the
\emph{metric complexity} $D_k\asymp2^{kd}$ of the model, i.e.\ a variance
$\asymp D_k/n$ with no logarithm. We make this precise on both routes flagged in
the lower-bound note.

Throughout, $\mathcal M_k$ is the linear space of functions piecewise polynomial
of degree $<\lceil s\rceil$ on the resolution-$2^{-k}$ partition of
Lemma~\ref{lem:jackson}, of dimension $D_k\asymp2^{kd}$, and
$\Pi_k:L^2(\mu_G)\to\mathcal M_k$ the $L^2(\mu_G)$-orthogonal projection. We use
the projection density estimator
\[
\widehat g_k:=\Pi_k\Big(\tfrac1n\sum_{i=1}^n\delta_{\pi(X_i)}\Big),\qquad
\widehat f_k:=\frac{T(\widehat g_k)}{\int T(\widehat g_k)\,d\mu_G},\qquad
T(u):=(u\vee\tfrac{c_0}{2})\wedge2C_1 ,
\]
the truncation $T$ keeping $\widehat f_k$ a density valued in $[c_0/2,2C_1]$ and,
since $T$ is a $1$-Lipschitz contraction onto an interval containing
$f^\star\in[c_0,C_1]$, changing the $L^2$ risk only by a universal factor.

\subsection{Route (i): known $s$, exact rate}

\begin{proposition}[Single-resolution projection estimator]\label{prop:knowns}
Fix $s$ and let $k=k_n$ satisfy $2^{k_n}\asymp n^{1/(2s+d)}$ (so $D_{k_n}\asymp
n^{d/(2s+d)}$). Then for every $f^\star\in\Sigma_{\mathrm{inv}}(s,L)$,
\[
\E_{P^\star}\big[h^2(P^\star,\widehat f_{k_n})\big]
\ \le\ C\Big(\underbrace{2^{-2k_n s}}_{\mathrm{bias}^2}+\underbrace{\tfrac{D_{k_n}}{n}}_{\mathrm{variance}}\Big)
\ \le\ C'\,n^{-2s/(2s+d)},
\]
the exact minimax rate of the companion lower bound; no logarithmic factor.
\end{proposition}

\begin{proof}
Since $\widehat f_k\in[c_0/2,2C_1]$ and $f^\star\in[c_0,C_1]$, Hellinger and
$L^2(\mu_G)$ are comparable, $h^2(P^\star,\widehat f_k)\le\frac1{c_0}\|\widehat
f_k-f^\star\|_2^2$, and the contraction $T$ together with renormalisation only
decreases the $L^2$ distance to $f^\star$ up to a universal factor; so it suffices
to bound $\E\|\widehat g_k-f^\star\|_2^2$. Orthogonality of the projection splits it
into bias and variance,
\[
\E\|\widehat g_k-f^\star\|_2^2
=\underbrace{\|f^\star-\Pi_kf^\star\|_2^2}_{\mathrm{bias}^2}
+\underbrace{\E\|\widehat g_k-\Pi_kf^\star\|_2^2}_{\mathrm{variance}} .
\]
\emph{Bias.} The projection is the $L^2$-best element of $\mathcal M_k$, hence at
least as good as the interpolant $q_k$ of Lemma~\ref{lem:jackson}:
$\|f^\star-\Pi_kf^\star\|_2\le\|f^\star-q_k\|_2\le CL\,2^{-k_n s}$.
\emph{Variance.} Let $\{e_1,\dots,e_{D_k}\}$ be an $L^2(\mu_G)$-orthonormal basis of
$\mathcal M_k$. Then $\widehat g_k-\Pi_kf^\star=\sum_\ell\big(\tfrac1n\sum_i
e_\ell(\pi(X_i))-\E e_\ell\big)\,e_\ell$, and by independence
\[
\E\|\widehat g_k-\Pi_kf^\star\|_2^2
=\frac1n\sum_{\ell=1}^{D_k}\operatorname{Var}\!\big(e_\ell(\pi(X))\big)
\le\frac1n\int\!\Big(\sum_{\ell}e_\ell^2\Big)f^\star\,d\mu_G
\le\frac{C_1}{n}\sum_{\ell}\|e_\ell\|_2^2=\frac{C_1D_k}{n} .
\]
Balancing $2^{-2k_n s}\asymp D_{k_n}/n=2^{k_n d}/n$, i.e.\ $2^{k_n(2s+d)}\asymp n$,
yields $C'n^{-2s/(2s+d)}$. The estimator uses $s$ only through $k_n$.
\end{proof}

\subsection{Route (ii): adaptive, exact rate}

\begin{theorem}[Penalised projection estimator]\label{thm:adaptive}
Let $\widehat k\in\argmin_k\big\{-\|\widehat g_k\|_2^2+\mathrm{pen}(k)\big\}$ with
penalty $\mathrm{pen}(k)=\kappa\,D_k/n$ for a large enough universal constant
$\kappa$, and $\widehat f:=\widehat f_{\widehat k}$. Then for every
$s\in(0,s_{\max}]$ and every $f^\star\in\Sigma_{\mathrm{inv}}(s,L)$,
\[
\E_{P^\star}\big[h^2(P^\star,\widehat f)\big]
\ \le\ C\,\inf_{k}\Big\{\|f^\star-\Pi_kf^\star\|_2^2+\frac{D_k}{n}\Big\}
\ \le\ C'\,n^{-2s/(2s+d)},
\]
\emph{adaptively in $s$}, with no logarithmic factor.
\end{theorem}

\begin{proof}
The criterion is the least-squares contrast $\gamma_n(g)=\|g\|_2^2-\frac2n\sum_i
g(\pi(X_i))$, minimised over $\mathcal M_k$ by $\widehat g_k=\Pi_k(\text{empirical})$
with $\gamma_n(\widehat g_k)=-\|\widehat g_k\|_2^2$; thus $\widehat k$ is the
penalised model-selection rule. The collection $\{\mathcal M_k\}_{k\ge0}$ is nested
with geometrically growing dimensions $D_k\asymp2^{kd}$, one model per dimension.
The Birg\'e--Massart penalised-projection oracle inequality for density estimation
\cite{birgemassart98,massart} applies: whenever the weights $\{x_k\}$ satisfy
$\Sigma:=\sum_k e^{-x_k}<\infty$, the selected estimator obeys
\[
\E\|\widehat g_{\widehat k}-f^\star\|_2^2
\le C\inf_k\Big\{\|f^\star-\Pi_kf^\star\|_2^2+\frac{D_k+x_k}{n}\Big\}+\frac{C\Sigma}{n}.
\]
Here is the decisive point. Because $D_k\asymp2^{kd}$ grows \emph{geometrically},
the choice $x_k=D_k$ already gives $\Sigma=\sum_k e^{-D_k}<\infty$ with $x_k\asymp
D_k$; the adaptation term $x_k/n$ is then of the \emph{same} order $D_k/n$ as the
variance and contributes no extra factor. (Contrast the MDL code of
Lemma~\ref{lem:code}, whose per-model weight is the codelength $L(k)\asymp
D_k\log n$, inflating the penalty -- and hence the rate -- by $\log n$.) Thus
$\mathrm{pen}(k)=\kappa D_k/n$ is admissible and
\[
\E\|\widehat g_{\widehat k}-f^\star\|_2^2\le C\inf_k\big\{\|f^\star-\Pi_kf^\star\|_2^2+D_k/n\big\}.
\]
Passing to $\widehat f$ as in Proposition~\ref{prop:knowns} (Hellinger
$\lesssim L^2$ under the two-sided bound) and inserting the Jackson bias
$\|f^\star-\Pi_kf^\star\|_2\le CL2^{-ks}$, the infimum is the balance of
Proposition~\ref{prop:knowns}, now realised by the data-driven $\widehat k$
\emph{without knowledge of $s$}: $\E\,h^2\le C'n^{-2s/(2s+d)}$.
\end{proof}

\begin{remark}[Honest status of the logarithmic factor: resolved]
The $(\log n)^{2s/(2s+d)}$ factor of Theorem~\ref{thm:upper} is the
two-stage/adaptation cost of \emph{description-length coding}
(Lemma~\ref{lem:code}), not an intrinsic feature of the problem: it is absent from
the lower bound, the exact minimax rate. Proposition~\ref{prop:knowns} (known $s$)
and Theorem~\ref{thm:adaptive} (adaptive) attain $n^{-2s/(2s+d)}$ \emph{exactly},
removing the log on both counts. The structural reason is that a projection /
penalised-projection estimator pays the metric complexity $D_k/n$ of the model,
whereas the MDL code pays an additional $\log(1/\delta)\asymp\log n$ per
coefficient to quantise the continuum. The plan may therefore state the
\emph{exact} minimax rate $n^{-2s/(2s+d)}$ for the orbital problem, with the MDL
estimator itself being near-minimax (up to $\log$) and the projection estimators
exact.
\end{remark}

\begin{remark}[What is proved vs.\ cited]
Fully proved here: the tail-free partition bound (Lemma~\ref{lem:partition}), the
twin Jackson estimate with the stratified-quotient resolution
(Lemma~\ref{lem:jackson}), the discretisation/description-length accounting
(Lemma~\ref{lem:code}), the balance (Theorem~\ref{thm:upper}), and the
log-removing bias/variance decomposition of the projection estimator
(Proposition~\ref{prop:knowns}). Cited as classical: the passage from
Lemma~\ref{lem:partition} to the resolvability oracle inequality (Barron--Cover
\cite{barroncover}, Li \cite{li}) and the penalised-projection oracle inequality
(Birg\'e--Massart \cite{birgemassart98,massart}). Both upper-side gaps flagged in
the lower-bound note are thereby closed, and the logarithmic factor between the
upper and lower bounds is removed.
\end{remark}

\begin{thebibliography}{9}
\bibitem{barroncover} A.~Barron, T.~Cover, ``Minimum complexity density
estimation,'' \emph{IEEE Trans.\ Inform.\ Theory} 37 (1991) 1034--1054.
\bibitem{li} J.~Q.~Li, \emph{Estimation of Mixture Models}, Ph.D.\ thesis, Yale
University, 1999.
\bibitem{tsybakov} A.~B.~Tsybakov, \emph{Introduction to Nonparametric
Estimation}, Springer, 2009.
\bibitem{wongshen} W.~H.~Wong, X.~Shen, ``Probability inequalities for likelihood
ratios and convergence rates of sieve MLEs,'' \emph{Ann.\ Statist.} 23 (1995)
339--362.
\bibitem{birgemassart98} L.~Birg\'e, P.~Massart, ``Minimum contrast estimators on
sieves: exponential bounds and rates of convergence,'' \emph{Bernoulli} 4 (1998)
329--375.
\bibitem{massart} P.~Massart, \emph{Concentration Inequalities and Model
Selection}, Lecture Notes in Math.\ 1896, Springer, 2007.
\end{thebibliography}

\end{document}
